% Luc Destombe --- licence CC BY-NC-SA 4.0
% https://creativecommons.org/licenses/by-nc-sa/4.0/deed.fr
% Fichier autonome : compiler avec pdflatex, deux fois (signets, nombre de pages).
\documentclass[a4paper,10pt]{article}
\makeatletter
\newif\iflycee@niveaux
\lycee@niveauxtrue
\newif\iflycee@palatino
\lycee@palatinotrue

\RequirePackage[utf8]{inputenc}
\RequirePackage[T1]{fontenc}
\RequirePackage[french]{babel}
\RequirePackage{etoolbox}

\newcommand{\ShorthandsOff}{\@ifundefined{shorthandoff}{}{\shorthandoff{;:!?}}}
\newcommand{\ShorthandsOn}{\@ifundefined{shorthandon}{}{\shorthandon{;:!?}}}
\ShorthandsOff
\AtBeginDocument{\ShorthandsOn}
\AtBeginEnvironment{tikzpicture}{\ShorthandsOff}

\RequirePackage{geometry}
\RequirePackage{fancyhdr}
\RequirePackage{lastpage}
\RequirePackage{multicol}
\RequirePackage{array}
\RequirePackage{enumitem}
\RequirePackage{xcolor}
\RequirePackage{needspace}

\RequirePackage{amsmath,amssymb}
\RequirePackage{mathpazo}
\DeclareMathSizes{10.95}{10.95}{8}{5.5}
\begingroup\catcode`\_=\active \catcode`\^=\active
\gdef_#1{\sb{\scriptscriptstyle #1}}
\gdef^#1{\sp{\scriptscriptstyle\lycee@moinsepais #1}}
\endgroup
\newcommand\lycee@moins{\mathbin{%
  \setbox\z@\hbox{$\m@th\scriptscriptstyle\lycee@moinsorig$}%
  \dimen@=.55\wd\z@ \advance\dimen@-.5pt
  \hbox to.75\wd\z@{\hss$\m@th\scriptscriptstyle\vcenter{\hbox{%
    \tikz\draw[line width=.5pt,line cap=round](0,0)--(\the\dimen@,0);}}$\hss}}}
\begingroup\lccode`\~=`\- \lowercase{\endgroup\let~}\lycee@moins
\newcommand\lycee@moinsepais{\mathcode`\-="8000 }
\AtBeginDocument{\mathchardef\lycee@moinsorig=\mathcode`\-\relax}
\AtBeginDocument{\catcode`\_=12 \mathcode`\_="8000
  \catcode`\^=12 \mathcode`\^="8000 }
\newcommand\lycee@exposants{%
  \fontdimen13\textfont2=.48\fontdimen6\textfont2
  \fontdimen14\textfont2=.48\fontdimen6\textfont2
  \fontdimen15\textfont2=.48\fontdimen6\textfont2
  \fontdimen13\scriptfont2=.48\fontdimen6\scriptfont2
  \fontdimen14\scriptfont2=.48\fontdimen6\scriptfont2
  \fontdimen15\scriptfont2=.48\fontdimen6\scriptfont2}
\every@math@size\expandafter{\the\every@math@size\lycee@exposants}
\newlength{\retraitcalculs}
\setlength{\retraitcalculs}{2em}

\RequirePackage{tikz}
\usetikzlibrary{arrows.meta,calc,babel,shapes.misc}
\RequirePackage[most]{tcolorbox}
\tcbuselibrary{skins,breakable}

\definecolor{coulDef}{HTML}{1F4E79}
\definecolor{coulProp}{HTML}{7030A0}
\definecolor{coulRem}{HTML}{C55A11}
\definecolor{coulEx}{HTML}{2E7D32}
\definecolor{coulExo}{HTML}{B03A2E}
\definecolor{coulPart}{HTML}{1F4E79}
\definecolor{coulMeth}{HTML}{1B6E4F}
\definecolor{coulCourbe}{HTML}{0B5ED7}
\definecolor{coulCourbeB}{HTML}{C00000}
\definecolor{coulCourbeC}{HTML}{2E7D32}
\definecolor{coulGrille}{HTML}{8C8C8C}
\definecolor{coulRep}{HTML}{1B6E4F}
\definecolor{coulBar}{HTML}{2E7D32}

\newcommand{\R}{\mathbb{R}}
\newcommand{\Rs}{\mathbb{R}^{*}}
\renewcommand{\le}{\leqslant}
\renewcommand{\ge}{\geqslant}
\newcommand{\vect}[1]{\overrightarrow{#1}}
\newcommand{\norme}[1]{\left\lVert #1 \right\rVert}
\newcommand{\intervalle}[2]{\left[#1\,;#2\right]}
\RequirePackage{cancel}
\renewcommand{\CancelColor}{\color{coulExo}}
\newcommand{\fois}[1]{\times{\color{coulExo}#1}}
\newcommand{\barre}[1]{\cancel{#1}}

\tikzset{grille/.style={coulGrille,line width=0.4pt}}
\newcommand{\quadrillage}[4]{\draw[grille,step=1] (#1,#2) grid (#3,#4);}
\tikzset{
  croix/.style={cross out,draw,line width=0.6pt,inner sep=0pt,minimum size=4pt},
  croixdroite/.style={croix,rotate=45,line width=0.8pt},
}
\newcommand{\pt}[3]{\path (#1) node[croix]{} node[#2,font=\small,inner sep=2.5pt]{$#3$};}

\newcommand{\lycee@repere}[4]{%
  \draw[grille,step=1] (#1,#2) grid (#3,#4);
  \draw[-{Stealth[length=2mm]},thick] (#1,0)--({#3+0.4},0) node[below left=-1pt]{$x$};
  \draw[-{Stealth[length=2mm]},thick] (0,#2)--(0,{#4+0.4}) node[below left=-1pt]{$y$};
  \node[below left,font=\scriptsize,inner sep=1.5pt] at (0,0) {$0$};}
\newcommand{\lycee@gradx}[1]{\foreach \k/\l in {#1}{%
  \draw (\k,0.07)--(\k,-0.07) node[below,font=\scriptsize,inner sep=1.5pt]{$\l$};}}
\newcommand{\lycee@grady}[1]{\foreach \k/\l in {#1}{%
  \draw (0.07,\k)--(-0.07,\k) node[left,font=\scriptsize,inner sep=1.5pt]{$\l$};}}
\newcommand{\lycee@intff}[2]{\left[#1\,;#2\right]}
\newcommand{\lycee@intfo}[2]{\left[#1\,;#2\right[}
\newcommand{\lycee@intof}[2]{\left]#1\,;#2\right]}
\newcommand{\lycee@intoo}[2]{\left]#1\,;#2\right[}
\newcommand{\lycee@ens}[1]{\left\{#1\right\}}
\AtBeginDocument{%
  \providecommand{\repere}{\lycee@repere}%
  \providecommand{\gradx}{\lycee@gradx}%
  \providecommand{\grady}{\lycee@grady}%
  \providecommand{\Cf}{\mathcal{C}\sb{\scriptscriptstyle f}}%
  \providecommand{\Cg}{\mathcal{C}\sb{\scriptscriptstyle g}}%
  \providecommand{\intff}{\lycee@intff}%
  \providecommand{\intfo}{\lycee@intfo}%
  \providecommand{\intof}{\lycee@intof}%
  \providecommand{\intoo}{\lycee@intoo}%
  \providecommand{\ens}{\lycee@ens}}

\newlist{questions}{enumerate}{3}
\setlist[questions,1]{label=\arabic*\textdegree),leftmargin=*,itemsep=2pt,topsep=3pt}
\setlist[questions,2]{label=\alph*),leftmargin=*,itemsep=1pt,topsep=2pt}
\setlist[questions,3]{label=\roman*),leftmargin=*,itemsep=1pt,topsep=1pt}
\setlist[itemize]{label=\textbullet,leftmargin=*,itemsep=2pt,topsep=3pt}

\newcommand{\besoinplace}[1]{\par\penalty\@M\begingroup
  \dimen@=#1\relax\dimen@ii=\pagegoal\advance\dimen@ii-\pagetotal
  \ifdim\dimen@>\dimen@ii\ifdim\dimen@ii>\z@\vfil\fi\break\fi\endgroup}

\newcommand{\lycee@pied}{\thepage/\pageref{LastPage}}
\newcommand{\entete}[2]{%
  \pagestyle{fancy}\fancyhf{}%
  \fancyhead[L]{\footnotesize\color{coulPart}\textbf{#1}}%
  \fancyhead[R]{\footnotesize\color{coulPart}\textbf{#2}}%
  \fancyfoot[C]{\footnotesize\lycee@pied}%
  \renewcommand{\headrulewidth}{0.4pt}%
  \renewcommand{\headrule}{\hbox to\headwidth{%
    \color{coulPart!40}\leaders\hrule height \headrulewidth\hfill}}}

\RequirePackage{ccicons}
\newcommand{\lycee@licence}{{\scriptsize\color{gray}%
  \copyright\ Luc Destombe --- \ccbyncsa\ CC BY-NC-SA 4.0}}
\newif\iflycee@licence \lycee@licencetrue
\newcommand{\sanslicence}{\lycee@licencefalse}
\AtBeginDocument{\iflycee@licence\AddToHookNext{shipout/foreground}{%
  \put(0,0){\rlap{\kern\dimexpr1in+\hoffset+\oddsidemargin+\textwidth\relax
    \llap{\raisebox{-\dimexpr1in+\voffset+\topmargin+\headheight+\headsep
      +\textheight+\footskip\relax}{\lycee@licence}}}}}\fi}

\newcommand{\formulesserrees}{%
  \setlength{\abovedisplayskip}{3pt}\setlength{\belowdisplayskip}{3pt}%
  \setlength{\abovedisplayshortskip}{2pt}\setlength{\belowdisplayshortskip}{3pt}}

\setlength{\parindent}{0pt}

\RequirePackage[implicit=false,hidelinks,pdfencoding=auto,psdextra,bookmarksopen,
  bookmarksopenlevel=1,pdfcreator={},pdfproducer={}]{hyperref}
\RequirePackage{bookmark}
\hypersetup{pdfauthor={Luc Destombe},pdfkeywords={CC BY-NC-SA 4.0}}
\newcounter{lycee@signet}
\newcommand{\signet}[3][]{\stepcounter{lycee@signet}%
  \edef\lycee@dest{\if\relax\detokenize{#1}\relax
    signet.\arabic{lycee@signet}\else#1\fi}%
  \hypertarget{\lycee@dest}{}\bookmark[level=#2,dest=\lycee@dest]{#3}}
\pdfstringdefDisableCommands{%
  \def\R{R}\def\vect#1{#1}\def\Cf{Cf}\def\Cg{Cg}\def\fois#1{×#1}%
  \def\barre#1{#1}\def\intervalle#1#2{[#1;#2]}\def\intff#1#2{[#1;#2]}%
  \def\textsuperscript#1{#1}\def\dfrac#1#2{#1/#2}\def\frac#1#2{#1/#2}%
  \def\vec#1{#1⃗}}%
\RequirePackage[official]{eurosym}

\geometry{top=0.8cm,bottom=1.3cm,left=1cm,right=1cm,footskip=0.6cm}

\pagestyle{fancy}
\fancyhf{}
\renewcommand{\headrulewidth}{0pt}
\fancyfoot[C]{\footnotesize Page \thepage{} / \pageref{LastPage}}

\newcommand{\titrefiche}[2]{%
  \begin{center}
    \Large\textbf{#1}\\[2pt]
    \large\textbf{#2}\\[2pt]
  \end{center}
  \vspace{0.10cm}}

\setlength{\columnseprule}{0.4pt}
\setlength{\columnsep}{15pt}
\renewcommand{\columnseprulecolor}{\color{black!45}}
\setlength{\multicolsep}{2pt}

\newcounter{partiefiche}
\renewcommand{\thepartiefiche}{\Roman{partiefiche}}
\newif\ifapresbandeau
\newcommand{\lycee@bandeau}[1]{%
  \par\penalty-600
  \vspace{0.08cm}%
  \noindent\colorbox{coulPart}{%
    \parbox{\dimexpr\linewidth-2\fboxsep\relax}{%
      \color{white}\bfseries\raggedright\strut#1\strut}}%
  \nopagebreak\par\nopagebreak
  \vspace{0.06cm}\nopagebreak
  \apresbandeautrue}
\newcommand{\partiefiche}[1]{%
  \refstepcounter{partiefiche}%
  \lycee@bandeau{\signet{1}{\thepartiefiche{} --- #1}\thepartiefiche{} --- #1}}
\newcommand{\partiefichesansnum}[1]{\lycee@bandeau{\signet{1}{#1}#1}}

\newcounter{exo}
\newlength{\espaceexo}\setlength{\espaceexo}{7pt}
\newlength{\espacetitre}\setlength{\espacetitre}{2pt}
\newcommand{\exo}[1][\relax]{%
  \par
  \ifapresbandeau\apresbandeaufalse\else\penalty-150\addvspace{\espaceexo}\fi
  \refstepcounter{exo}%
  \noindent\signet[exercice-\arabic{exo}]{2}{Exercice \arabic{exo}}%
  \textbf{\color{coulExo}\underline{Exercice \theexo}}%
  \ifx\relax#1\relax\else\textbf{ : #1}\fi
  \par\nopagebreak\vspace{\espacetitre}\nopagebreak
  \ignorespaces}

\setlist[enumerate]{nosep,leftmargin=*,itemsep=2pt,topsep=2pt}
\setlist[itemize]{nosep,leftmargin=*,topsep=2pt}
\newcommand{\choix}[3]{\par\hspace*{1em}%
  \makebox[0.3\linewidth][l]{a) #1}\makebox[0.3\linewidth][l]{b) #2}c) #3}
\newcommand{\deux}[2]{\par\hspace*{0.6em}%
  \makebox[0.48\linewidth][l]{#1}#2}
\newcommand{\trois}[3]{\par\hspace*{0.6em}%
  \makebox[0.32\linewidth][l]{#1}\makebox[0.32\linewidth][l]{#2}#3}

  \renewcommand{\exo}[1][\relax]{%
    \ifapresbandeau\apresbandeaufalse\else\par\penalty-150\fi
    \refstepcounter{exo}%
    \vspace{0.05cm}%
    \noindent\signet[exercice-\arabic{exo}]{2}{Exercice \arabic{exo}}%
    \textbf{\color{coulExo}\underline{Exercice \theexo}}%
    \ifx\relax#1\relax\else\textbf{ : #1}\fi%
    \par\nopagebreak
    \ignorespaces}
  \newcommand{\rappel}[1]{%
    \par\vspace{0.06cm}\nopagebreak
    \noindent\fcolorbox{black!35}{black!5}{%
      \parbox{\dimexpr\linewidth-2\fboxsep-2\fboxrule\relax}{%
        \small\textbf{Rappel.} #1}}%
    \par\vspace{0.06cm}\nopagebreak}
  \setlist[enumerate]{nosep,leftmargin=*}
  \setlist[itemize]{nosep,leftmargin=*}
  \setlength{\multicolsep}{3pt plus 1pt minus 1pt}
  \setlength{\premulticols}{10pt}
  \setlength{\postmulticols}{10pt}
  \newlist{expr}{enumerate}{1}
  \setlist[expr]{label={\Alph*\ $=$},nosep,leftmargin=*,itemsep=0pt}

\renewenvironment{center}{\par\addvspace{3pt}\centering}{\par\addvspace{3pt}}
\formulesserrees
\AtBeginDocument{\formulesserrees}
\clubpenalty=10000
\widowpenalty=10000
\displaywidowpenalty=10000
\brokenpenalty=10000
\setlength{\parskip}{0pt}

\RequirePackage{ragged2e}
\setlength{\RaggedRightRightskip}{0pt plus 1fil}
\AtBeginDocument{\RaggedRight\binoppenalty=10000 \relpenalty=10000 }
\g@addto@macro\@parboxrestore{\RaggedRight}
\makeatother
\usepackage{tkz-tab}
\geometry{top=0.8cm,bottom=1.3cm,left=1cm,right=1cm,footskip=0.7cm,headheight=12pt}

\begin{document}
\titrefiche{1\textsuperscript{re} spécialité --- Fonctions}{Fiche d'exercices du chapitre}

\begin{multicols}{2}

\partiefiche{Images et antécédents}

\exo[Vocabulaire]
\begin{enumerate}[label=\arabic*),topsep=0pt]
  \item Écrire chaque phrase sous la forme d'une égalité.
    \begin{enumerate}[label=\alph*)]
      \item L'image de $3$ par la fonction $f$ est $-1$.
      \item Le nombre $-2$ est un antécédent de $6$ par la fonction $g$.
      \item Le point $B(0\,;5)$ est sur la courbe de la fonction $h$.
      \item Par la fonction $k$, le nombre $4$ a pour image $9$.
    \end{enumerate}
  \item On sait que $f(5)=-7$. Écrire trois phrases : la première avec le mot
        « image », la deuxième avec « antécédent », la troisième avec « courbe ».
\end{enumerate}

\exo[Calculer des images]
Soit $f$ la fonction définie sur $\R$ par $f(x)=2x^{2}-5x$.
\begin{enumerate}[label=\arabic*)]
  \item Calculer l'image par $f$ de $3$ ; de $-2$ ; de $\frac{1}{2}$ ; de $-1{,}5$ ; de
        $\sqrt{3}$.
  \item Le nombre $-1$ est-il un antécédent de $7$ ? Le nombre $2$ a-t-il pour image
        $2$ ?
\end{enumerate}

\exo[Calculer des antécédents]
Pour chaque fonction, déterminer les antécédents du nombre indiqué.
\begin{enumerate}[label=\alph*)]
  \item $f(x)=x-7$ ; antécédents de $2$.
  \item $g(x)=x^{2}$ ; antécédents de $49$.
  \item $h(x)=5x+3$ ; antécédents de $0$.
  \item $k(x)=(x-1)(x+4)$ ; antécédents de $0$.
  \item $m(x)=x^{2}+3$ ; antécédents de $2$.
\end{enumerate}

\exo[Lire des images et des antécédents]
La courbe $\Cg$ d'une fonction $g$ définie sur $\intff{-3}{5}$ est tracée ci-dessous.
\begin{center}
\begin{tikzpicture}[x=0.5cm,y=0.5cm]
  \repere{-4}{-3}{6}{5}
  \gradx{-3,-2,-1,1,2,3,4,5} \grady{-2,-1,1,2,3,4}
  \draw[coulCourbe,line width=1.1pt] (-3,4)
    .. controls (-2.667,3) and (-2.333,1.833) .. (-2,1)
    .. controls (-1.667,0.167) and (-1.333,-0.5) .. (-1,-1)
    .. controls (-0.667,-1.5) and (-0.333,-2) .. (0,-2)
    .. controls (0.333,-2) and (0.667,-1.5) .. (1,-1)
    .. controls (1.333,-0.5) and (1.667,0.5) .. (2,1)
    .. controls (2.333,1.5) and (2.667,2) .. (3,2)
    .. controls (3.333,2) and (3.667,1.5) .. (4,1)
    .. controls (4.333,0.5) and (4.667,-0.333) .. (5,-1);
  \path[coulCourbe] (-3,4) node[croix]{} (5,-1) node[croix]{};
  \node[coulCourbe,font=\footnotesize] at (-2.4,4.3) {$\Cg$};
\end{tikzpicture}
\end{center}
\begin{enumerate}[label=\arabic*)]
  \item Lire les images de $-2$ ; $0$ ; $3$ ; $5$.
  \item Lire les antécédents de $1$ ; $-1$ ; $-2$ ; $5$.
  \item Combien de solutions l'équation $g(x)=0$ a-t-elle ?
  \item Compléter : $g(\ldots)=4$.
\end{enumerate}

\exo[Points d'une courbe]
Soit $h$ la fonction définie sur $\R$ par $h(x)=x^{2}+2x-8$, de courbe $\mathcal{C}_h$.
\begin{enumerate}[label=\arabic*)]
  \item Les points $A(1\,;-5)$ et $B(-2\,;0)$ sont-ils sur $\mathcal{C}_h$ ?
  \item Calculer les coordonnées du point où $\mathcal{C}_h$ coupe l'axe des
        ordonnées.
  \item Développer $(x-2)(x+4)$, puis en déduire les points où $\mathcal{C}_h$ coupe
        l'axe des abscisses.
\end{enumerate}

\partiefiche{Ensemble de définition}

\exo[Un nombre sans image]
Pour chaque fonction $f$, $g$, $h$, l'un des nombres $-2$ ; $0$ ; $5$ ; $10$ n'a pas
d'image. Lequel ?
\[
  f(x)=\frac{3}{x} \qquad g(x)=\frac{x+3}{x-5} \qquad h(x)=\sqrt{x+1}
\]

\columnbreak
\exo[Ensembles de définition]
Déterminer l'ensemble de définition de chaque fonction.
\deux{a) $f(x)=\dfrac{2}{x+6}$}{b) $g(x)=\dfrac{5}{4x-12}$}\\[3pt]
\deux{c) $h(x)=\dfrac{x}{x^{2}+5}$}{d) $k(x)=\dfrac{1}{x^{2}-16}$}\\[3pt]
\deux{e) $m(x)=\sqrt{12-3x}$}{f) $p(x)=\dfrac{1}{\sqrt{x-1}}$}

\exo[Avec une valeur interdite]
Soit $f$ la fonction définie par $f(x)=\dfrac{6}{x-1}$.
\begin{enumerate}[label=\arabic*)]
  \item Déterminer l'ensemble de définition de $f$.
  \item Calculer $f(0)$ ; $f(4)$ ; $f(-2)$.
  \item Déterminer les antécédents de $3$ ; de $-6$ ; de $0$.
\end{enumerate}

\exo[Avec une racine carrée]
Soit $g$ la fonction définie par $g(x)=\sqrt{4x+8}$.
\begin{enumerate}[label=\arabic*)]
  \item Déterminer l'ensemble de définition de $g$.
  \item Calculer, si possible, $g(2)$ ; $g(-1)$ ; $g(-3)$.
  \item Déterminer les antécédents de $6$ ; de $0$ ; de $-1$.
\end{enumerate}

\partiefiche{Comparaison de deux fonctions}

\exo[Comparer sur un graphique]
Les courbes de deux fonctions $f$ et $g$ définies sur $\intff{-3}{5}$ sont tracées
ci-dessous.
\begin{center}
\begin{tikzpicture}[x=0.5cm,y=0.5cm]
  \repere{-4}{-3}{6}{4}
  \gradx{-3,-2,-1,1,2,3,4,5} \grady{-2,-1,1,2,3}
  \draw[coulCourbe,line width=1.1pt] (-3,3)
    .. controls (-2.667,2.333) and (-2.333,1.667) .. (-2,1)
    .. controls (-1.667,0.333) and (-1.333,-0.5) .. (-1,-1)
    .. controls (-0.667,-1.5) and (-0.333,-2) .. (0,-2)
    .. controls (0.333,-2) and (0.667,-0.667) .. (1,0)
    .. controls (1.333,0.667) and (1.667,1.5) .. (2,2)
    .. controls (2.333,2.5) and (2.667,3) .. (3,3)
    .. controls (3.333,3) and (3.667,2.5) .. (4,2)
    .. controls (4.333,1.5) and (4.667,0.667) .. (5,0);
  \draw[coulCourbeC,line width=1.1pt] (-3,-1)
    .. controls (-2.667,-0.333) and (-2.333,0.5) .. (-2,1)
    .. controls (-1.667,1.5) and (-1.333,2) .. (-1,2)
    .. controls (-0.667,2) and (-0.333,1.333) .. (0,1)
    .. controls (0.333,0.667) and (0.667,0.333) .. (1,0)
    .. controls (1.333,-0.333) and (1.667,-1) .. (2,-1)
    .. controls (2.333,-1) and (2.667,-0.5) .. (3,0)
    .. controls (3.333,0.5) and (3.667,1.5) .. (4,2)
    .. controls (4.333,2.5) and (4.667,2.667) .. (5,3);
  \node[coulCourbe,font=\footnotesize] at (2.6,3.6) {$\Cf$};
  \node[coulCourbeC,font=\footnotesize] at (5.5,3.2) {$\Cg$};
\end{tikzpicture}
\end{center}
Résoudre graphiquement :
\begin{enumerate}[label=\alph*)]
  \item l'équation $f(x)=g(x)$ ;
  \item l'inéquation $f(x)>g(x)$ (réponse avec des intervalles) ;
  \item l'inéquation $g(x)\ge 2$.
\end{enumerate}

\exo[Signe d'une différence]
Pour chaque couple de fonctions définies sur $\R$, calculer $f(x)-g(x)$, étudier son
signe, puis en déduire la position relative de $\Cf$ et $\Cg$.
\begin{enumerate}[label=\alph*)]
  \item $f(x)=3x-4$ et $g(x)=x+2$
  \item $f(x)=x^{2}$ et $g(x)=2x-1$
  \item $f(x)=x^{2}-x$ et $g(x)=-x+4$
\end{enumerate}

\exo[Une parabole et une droite]
Soient $f$ et $g$ définies sur $\R$ par $f(x)=x^{2}-2x-2$ et $g(x)=x+2$.
\begin{enumerate}[label=\arabic*)]
  \item Dresser un tableau de valeurs de $f$ et de $g$ pour $x$ entier de $-2$ à $5$,
        puis tracer $\Cf$ et $\Cg$ dans un même repère.
  \item Résoudre graphiquement $f(x)=g(x)$, puis $f(x)\le g(x)$.
  \item Démontrer que $f(x)-g(x)=(x+1)(x-4)$.
  \item Dresser le tableau de signes de $f(x)-g(x)$.
  \item En déduire la position relative de $\Cf$ et $\Cg$ ; comparer avec le 2).
\end{enumerate}

\exo[Une fonction quotient]
Soit $f$ la fonction définie par $f(x)=\dfrac{2x+1}{x-3}$.
\begin{enumerate}[label=\arabic*)]
  \item Déterminer l'ensemble de définition $D$ de $f$.
  \item Calculer $f(0)$ ; $f(4)$ ; $f(2)$.
  \item Les nombres $1$ ; $0$ ; $2$ ont-ils des antécédents ? Si oui, lesquels ?
  \item Justifier que, pour tout $x$ de $D$, $f(x)=2+\dfrac{7}{x-3}$.
  \item En déduire la position de $\Cf$ par rapport à la droite d'équation $y=2$.
  \item Calculer $f(x)$ pour $x=-4$ ; $-1$ ; $5$ ; $10$ (au centième près), puis tracer
        $\Cf$ et la droite d'équation $y=2$.
\end{enumerate}

\partiefiche{Sens de variation}

\exo[Lecture graphique complète]
La courbe $\Cf$ d'une fonction $f$ définie sur $\intff{-4}{6}$ est tracée ci-dessous,
avec la droite $(D)$ d'équation $y=-x+1$.
\begin{center}
\begin{tikzpicture}[x=0.5cm,y=0.5cm]
  \repere{-5}{-4}{7}{5}
  \gradx{-4,-3,-2,-1,1,2,3,4,5,6} \grady{-3,-2,-1,1,2,3,4}
  \draw[coulCourbe,line width=1.1pt] (-4,0)
    .. controls (-3.667,0.333) and (-3.333,0.667) .. (-3,1)
    .. controls (-2.667,1.333) and (-2.333,2) .. (-2,2)
    .. controls (-1.667,2) and (-1.333,1.333) .. (-1,1)
    .. controls (-0.667,0.667) and (-0.333,0.5) .. (0,0)
    .. controls (0.333,-0.5) and (0.667,-1.5) .. (1,-2)
    .. controls (1.333,-2.5) and (1.667,-3) .. (2,-3)
    .. controls (2.333,-3) and (2.667,-2.5) .. (3,-2)
    .. controls (3.333,-1.5) and (3.667,-0.5) .. (4,0)
    .. controls (4.333,0.5) and (4.667,0.5) .. (5,1)
    .. controls (5.333,1.5) and (5.667,2.333) .. (6,3);
  \path[coulCourbe] (-4,0) node[croix]{} (6,3) node[croix]{};
  \draw[coulCourbeC,line width=1.1pt] (-3,4) -- (5,-4);
  \node[coulCourbe,font=\footnotesize] at (5.4,3.4) {$\Cf$};
  \node[coulCourbeC,font=\footnotesize] at (-2.4,4.4) {$(D)$};
\end{tikzpicture}
\end{center}
\begin{enumerate}[label=\arabic*)]
  \item Lire les antécédents de $0$ ; $1$ ; $-2$ ; $4$.
  \item Donner le maximum et le minimum de $f$ sur $\intff{-4}{6}$, et en quelles
        valeurs ils sont atteints.
  \item Dresser le tableau de variations de $f$.
  \item Résoudre graphiquement $f(x)\ge 0$, puis $f(x)<1$.
  \item Compléter : si $-2\le x\le 2$, alors $\ldots\le f(x)\le\ldots$ ; si $2\le x\le 5$,
        alors $\ldots\le f(x)\le\ldots$
  \item Résoudre graphiquement $f(x)=-x+1$, puis $f(x)\le -x+1$.
\end{enumerate}

\exo[Du tableau de variations à la courbe]
Voici le tableau de variations d'une fonction $g$.
\begin{center}
\begin{tikzpicture}[scale=0.7]
  \tkzTabInit[lgt=1.4,espcl=2]{$x$/1,$g(x)$/1.6}{$-3$,$0$,$2$,$4$}
  \tkzTabVar{-/$-1$,+/$3$,-/$-2$,+/$1$}
\end{tikzpicture}
\end{center}
\begin{enumerate}[label=\arabic*)]
  \item Quel est l'ensemble de définition de $g$ ?
  \item Décrire les variations de $g$ par des phrases.
  \item Donner le maximum et le minimum de $g$, et où ils sont atteints.
  \item Tracer une courbe qui peut représenter $g$.
\end{enumerate}

\exo[Le bénéfice d'un atelier]
Un atelier assemble et vend des trottinettes électriques. Pour $x$ dizaines de
trottinettes vendues dans la semaine, avec $x\in\intff{1}{11}$, le bénéfice, en centaines
d'euros, est $B(x)=-x^{2}+12x-20$.
\begin{enumerate}[label=\arabic*)]
  \item Calculer $B(x)$ pour $x$ entier de $1$ à $11$.
  \item Tracer la courbe de $B$ (en abscisses $1$~cm pour une dizaine, en ordonnées
        $1$~cm pour $200$~€).
  \item Dresser le tableau de variations de $B$.
  \item Quel est le bénéfice maximal ? Pour combien de trottinettes ?
  \item Pour quelles ventes l'atelier ne perd-il pas d'argent ($B(x)\ge 0$) ?
\end{enumerate}

\exo[Sens de variation par le calcul]
Dans tout l'exercice, $a$ et $b$ sont deux réels tels que $a<b$.
\begin{enumerate}[label=\arabic*)]
  \item Soit $f(x)=4x-7$ sur $\R$. Calculer et factoriser $f(a)-f(b)$, en déduire son
        signe, puis le sens de variation de $f$.
  \item Même travail avec $g(x)=-3x+2$.
  \item Que peut-on conclure pour une fonction affine $x\mapsto mx+p$, selon le signe
        de $m$ ?
  \item Soit $u(x)=x^{2}+4x$ sur $\intfo{0}{+\infty}$. Démontrer que
        $u(a)-u(b)=(a-b)(a+b+4)$, puis conclure.
  \item Soit $v(x)=1-\dfrac{1}{x}$ sur $\intoo{0}{+\infty}$. Démontrer que
        $v(a)-v(b)=\dfrac{a-b}{ab}$, puis conclure.
\end{enumerate}

\partiefiche{Fonction carré}

\exo[Équations avec des carrés]
Résoudre dans $\R$ :
\trois{a) $x^{2}=36$}{b) $x^{2}=11$}{c) $2x^{2}-5=x^{2}+4$}
\deux{d) $(3x+1)^{2}=16$}{e) $x^{2}+9=0$}

\exo[Comparer et encadrer des carrés]
\begin{enumerate}[label=\arabic*),topsep=0pt]
  \item Comparer sans calculatrice, en citant le sens de variation utilisé :
        $(-4{,}2)^{2}$ et $(-3{,}9)^{2}$ ; puis $1{,}05^{2}$ et $1{,}5^{2}$.
  \item Encadrer $x^{2}$ sachant que $3\le x\le 6$ ; puis $-5\le x\le -2$ ; puis
        $-3\le x\le 1$.
\end{enumerate}

\exo[Deux paraboles]
Soient $f$ et $g$ définies sur $\R$ par $f(x)=2x^{2}+x+7$ et $g(x)=x^{2}-3x+3$.
\begin{enumerate}[label=\arabic*)]
  \item Calculer $f(x)-g(x)$, puis vérifier que $f(x)-g(x)=(x+2)^{2}$.
  \item En déduire la position relative de $\Cf$ et $\Cg$, et leur point commun.
\end{enumerate}

\partiefiche{Fonction cube}

\exo[Variations de la fonction cube]
Soient $a$ et $b$ deux réels tels que $a<b$.
\begin{enumerate}[label=\arabic*)]
  \item Développer $(a-b)\left(a^{2}+ab+b^{2}\right)$ et vérifier que l'on obtient
        $a^{3}-b^{3}$.
  \item On suppose $a<b\le 0$. Justifier que $a^{2}+ab+b^{2}>0$, puis en déduire le
        signe de $a^{3}-b^{3}$ et le sens de variation de la fonction cube sur
        $\intof{-\infty}{0}$.
  \item Même travail pour $0\le a<b$. Conclure sur $\R$.
\end{enumerate}

\exo[Calculs et équations]
\begin{enumerate}[label=\arabic*),topsep=0pt]
  \item Calculer le cube de $3$ ; de $-2$ ; de $\frac{1}{3}$ ; de $-0{,}2$.
  \item Résoudre dans $\R$ :
        \deux{a) $x^{3}=-125$}{b) $x^{3}=1$}
        \deux{c) $3x^{3}+24=0$}{d) $x^{3}=4x$}
\end{enumerate}

\exo[Encadrer des cubes]
\begin{enumerate}[label=\arabic*),topsep=0pt]
  \item Encadrer $x^{3}$ sachant que $-3\le x\le 2$.
  \item Comparer sans calculatrice $(-2{,}1)^{3}$ et $(-1{,}9)^{3}$.
  \item Une boîte cubique a une arête de $x$~cm. Son arête est comprise entre $2$ et
        $3$~cm : encadrer son volume. Quelle arête donne un volume de $64$~cm$^{3}$ ?
\end{enumerate}

\partiefiche{Fonction inverse}

\exo[Variations de la fonction inverse]
\begin{enumerate}[label=\arabic*),topsep=0pt]
  \item Soient $0<a<b$. Justifier que $\dfrac{1}{a}-\dfrac{1}{b}=\dfrac{b-a}{ab}$, puis
        que ce nombre est positif. En déduire le sens de variation de la fonction
        inverse sur $\intoo{0}{+\infty}$.
  \item Même travail pour $a<b<0$.
  \item Dresser le tableau de variations de la fonction inverse.
  \item Comparer $\dfrac{1}{-2}$ et $\dfrac{1}{2}$. La fonction inverse est-elle
        décroissante sur $\Rs$ ?
\end{enumerate}

\exo[Équations et encadrements]
\begin{enumerate}[label=\arabic*),topsep=0pt]
  \item Résoudre :
        \deux{a) $\dfrac{1}{x}=8$}{b) $\dfrac{1}{x}=-\dfrac{2}{3}$}
        \deux{c) $\dfrac{4}{x}=6$}{d) $\dfrac{5}{x}=\dfrac{x}{5}$}
  \item Encadrer $\dfrac{1}{x}$ sachant que $2\le x\le 8$ ; puis $-10\le x\le -5$.
  \item Comparer sans calculatrice $\dfrac{1}{0{,}3}$ et $\dfrac{1}{0{,}4}$.
\end{enumerate}

\exo[Équations et inéquations quotients]
\begin{enumerate}[label=\arabic*),topsep=0pt]
  \item Résoudre $\dfrac{x+3}{x-2}=\dfrac{3}{2}$, après avoir donné la valeur interdite.
  \item Résoudre en se ramenant à un quotient comparé à $0$, puis avec un tableau de
        signes :
        \deux{a) $\dfrac{1}{x}-3\le 0$}{b) $\dfrac{4}{x-1}\ge 2$}
\end{enumerate}

\partiefiche{Fonction racine carrée}

\exo[Calculer avec des racines carrées]
\begin{enumerate}[label=\arabic*),topsep=0pt]
  \item Écrire plus simplement :
        \deux{a) $\sqrt{18}-\sqrt{2}$}{b) $\left(\sqrt{5}-\sqrt{2}\right)\left(\sqrt{5}+\sqrt{2}\right)$}
        \deux{c) $\left(\sqrt{3}+\sqrt{6}\right)^{2}$}{}
  \item Soient $A=3-\sqrt{10}$ et $B=\sqrt{19-6\sqrt{10}}$. Calculer $A^{2}$ et $B^{2}$.
        A-t-on $A=B$ ?
  \item Justifier les égalités :
        \deux{a) $\dfrac{1}{\sqrt{3}}=\dfrac{\sqrt{3}}{3}$}{b) $\dfrac{1}{\sqrt{7}-2}=\dfrac{\sqrt{7}+2}{3}$}
        \deux{c) $\sqrt{50}-\sqrt{32}+\sqrt{2}=2\sqrt{2}$}{}
  \item Calculer $\left(\sqrt{6}+\sqrt{7}\right)\left(\sqrt{6}-\sqrt{7}\right)$ ; en
        déduire, sans calculatrice, que $\sqrt{6}<\sqrt{7}$.
\end{enumerate}

\exo[Équations avec des racines carrées]
Résoudre, en précisant d'abord les valeurs de $x$ possibles :
\deux{a) $\sqrt{x}=9$}{b) $\sqrt{x}=-3$}
\deux{c) $\sqrt{x-4}=5$}{d) $3\sqrt{2x+1}=15$}

\exo[Se ramener au second degré]
On considère l'équation $(E)$ : $\sqrt{x+3}=x+1$.
\begin{enumerate}[label=\arabic*)]
  \item À la calculatrice, conjecturer le nombre de solutions de $(E)$ et leur valeur.
  \item Démontrer que toute solution de $(E)$ vérifie $x+1\ge 0$ et $x^{2}+x-2=0$.
  \item Vérifier que $x^{2}+x-2=(x+2)(x-1)$, puis résoudre $(E)$.
  \item Même démarche pour $2\sqrt{x}=x-3$ : se ramener à $x^{2}-10x+9=0$, puis
        vérifier que $x^{2}-10x+9=(x-5)^{2}-4^{2}$ pour factoriser.
\end{enumerate}

\partiefiche{Fonction valeur absolue}

\exo[Calculs et distances]
\begin{enumerate}[label=\arabic*),topsep=0pt]
  \item Calculer :
        \trois{a) $|-12|$}{b) $|3-10|$}{c) $|-4|\times|-5|$}
        \deux{d) $|\sqrt{2}-2|$}{e) $|-6|-|2|$}
  \item Calculer $\sqrt{(-8)^{2}}$ et $\sqrt{(-0{,}5)^{2}}$. Que vaut $\sqrt{a^{2}}$
        pour un réel $a$ quelconque ?
  \item Calculer la distance entre $-6$ et $3$ ; entre $7$ et $-1$.
\end{enumerate}

\exo[Équations avec des valeurs absolues]
Résoudre dans $\R$ :
\deux{a) $|x|=9$}{b) $|x|=-3$}
\deux{c) $|x+2|=6$}{d) $|3x-1|=8$}

\exo[Étude de $x\mapsto|x+3|$]
Soit $f$ la fonction définie sur $\R$ par $f(x)=|x+3|$.
\begin{enumerate}[label=\arabic*)]
  \item Calculer $f(x)$ pour $x$ entier de $-6$ à $0$, puis tracer la courbe de $f$.
  \item Comment se déduit-elle de la courbe de la valeur absolue ?
  \item Dresser le tableau de variations de $f$ ; donner son minimum.
  \item Résoudre graphiquement, puis par le calcul :
        \deux{a) $f(x)=2$}{b) $f(x)\le 1$}
\end{enumerate}

\partiefiche{Fonctions affines}

\exo[Lire l'expression d'une fonction affine]
Donner l'expression de la fonction affine représentée par chaque droite $(d_{1})$,
$(d_{2})$, $(d_{3})$.
\begin{center}
\begin{tikzpicture}[x=0.5cm,y=0.5cm]
  \repere{-5}{-4}{5}{4}
  \gradx{-4,-3,-2,-1,1,2,3,4} \grady{-3,-2,-1,1,2,3}
  \begin{scope}
  \clip (-5,-4) rectangle (5,4);
  \draw[coulCourbe,line width=1.1pt] (-1.5,4) -- (2.5,-4);
  \draw[coulCourbeB,line width=1.1pt] (-5,3) -- (5,3);
  \draw[coulCourbeC,line width=1.1pt] (-3,-4) -- (5,{4/3});
  \end{scope}
  \node[coulCourbe,font=\footnotesize] at (-0.3,3.6) {$(d_{1})$};
  \node[coulCourbeB,font=\footnotesize] at (-4.1,3.5) {$(d_{2})$};
  \node[coulCourbeC,font=\footnotesize] at (4.2,0.4) {$(d_{3})$};
\end{tikzpicture}
\end{center}

\exo[Équation réduite et alignement]
On donne les points $A(-1\,;4)$ et $B(2\,;-2)$.
\begin{enumerate}[label=\arabic*)]
  \item Calculer le coefficient directeur de la droite $(AB)$.
  \item En déduire l'équation réduite de $(AB)$.
  \item Les points $C(5\,;-8)$ et $D(-3\,;9)$ sont-ils alignés avec $A$ et $B$ ?
\end{enumerate}

\exo[Intersection de deux droites]
Calculer les coordonnées du point d'intersection des droites $(d)$ et $(d')$.
\begin{enumerate}[label=\arabic*)]
  \item $(d) : y=3x-5$ et $(d') : y=-2x+5$.
  \item $(d)$ passe par $A(0\,;-1)$ et $B(3\,;5)$ ; $(d') : y=-x+5$.
\end{enumerate}

\partiefiche{Prolongement : vers la dérivation}

\exo[Calculer des taux de variation]
\begin{enumerate}[label=\arabic*),topsep=0pt]
  \item $f(x)=x^{2}+2x$ : taux de variation de $f$ entre $0$ et $3$, puis entre $-4$
        et $-1$.
  \item $g(x)=-2x+7$ : taux de variation de $g$ entre $-1$ et $4$, puis entre $2$ et
        $10$. Que remarque-t-on ?
  \item Démontrer que, pour $x\mapsto mx+p$, le taux de variation entre deux nombres
        distincts vaut toujours $m$. Interpréter sur la droite.
\end{enumerate}

\exo[Taux et sens de variation]
\begin{enumerate}[label=\arabic*),topsep=0pt]
  \item Soit $g(x)=x^{2}-4x$. Démontrer que, pour $a\ne b$, le taux de variation de $g$
        entre $a$ et $b$ vaut $a+b-4$. En déduire le sens de variation de $g$ sur
        $\intfo{2}{+\infty}$, puis sur $\intof{-\infty}{2}$.
  \item Soit $h(x)=\dfrac{1}{x}$. Démontrer que, pour $a$ et $b$ non nuls et distincts,
        le taux de variation de $h$ entre $a$ et $b$ vaut $-\dfrac{1}{ab}$. En déduire
        le sens de variation de $h$ sur $\intoo{0}{+\infty}$.
\end{enumerate}

\exo[Vitesse moyenne d'une chute]
On lâche une bille du haut d'une tour. Au bout de $t$ secondes, elle a parcouru
$d(t)=5t^{2}$ mètres. Sa vitesse moyenne entre deux instants est le taux de variation
de $d$ entre ces instants.
\begin{enumerate}[label=\arabic*)]
  \item Calculer la vitesse moyenne entre $1$ et $3$ secondes.
  \item Même question entre $2$ et $2{,}5$ s, puis entre $2$ et $2{,}1$ s.
  \item Quelle semble être la vitesse de la bille à l'instant $t=2$ ?
\end{enumerate}

\vfill

\end{multicols}
\end{document}
